\scanpage{31}\mspage{24 ter}
% @ top and middle, subgroup | text | le groupe d'autom.
le groupe d'autom. de $(C,t,b)$ est le sous-groupe de $W_{D_4}$ formé des autom de $(E(w)^*,\sigma,-)$ qui invarient $v'\in E(w)^*$, c'est donc aussi le groupe $W_{D_3}$ des autom de $\bigl(E(w)^*-\{\unsure{w'|v'|u'}+ \sigma v'\},\sigma\bigr)$ qui correspondent à un nb pair de chgts de signes. Quel est le syst. de racines de type $D_3$ associé (formé de 12 racines) ? On trouve \deleted{sans doute le} syst. des racines de la forme $\alpha-\beta$, où $\alpha,\unsure{\beta'|\beta}\in E(w)^*\smallsetminus\{v',\sigma v'\}$ \inserted{sont non liés -- il ne dépend que de l'ens $\{s,t\}$ des}
\margin{deux triangles $s,t$ tels que $s\cap t\ne\varnothing$ $s\ne t$ (donc $\operatorname{Card}s\cap t=1$) ce sont aussi les racines orth. à $\{u,v,w,u',v'\}=s\cup t$ i.e. orth. à la fois $s$ et à $t$}% slanted

% @ middle, base with marked elements | text | La donnée de (t,b)
La donnée de $(t,b)$ équivaut aussi à celle d'une \uline{base} $b$, avec deux éléments fixés $\unsure{\zeta|\xi}_1$ et $\zeta_2$ (le triangle $t$ étant alors $(\underset{u}{\underbrace{\zeta_1}},\underset{v}{\underbrace{\zeta'_2}},\zeta_{12})$) -- le gpe des autom de $(C,t,b)$ est donc aussi $\unsure{\mathfrak S|\mathscr S}_4\simeq W_{A_3}$ -- on est dans le cas d'un ``isom. exceptionnel'' $D_3\simeq A_3$. Le nb de ces syst. $t,b$ est $72.6.5=2160=27.10.8$ --
% archive sheet number: 30
